\section{Limitations and conclusion}
\label{sec:conclusion}

\paragraph{Limitations.} We report results for one model; results on gpt-oss~\citep{gptoss} are coming soon. We
evaluate GRAFT at 3 bits per weight. Its tile-level advantage grows with the rate (\cref{fig:rd}), but we do not
evaluate higher rates end to end here. Model-level metrics use 400 windows from two text corpora, and we do not report
capability benchmarks. We do not measure throughput or latency, and the reference-kernel evaluation decodes each
layer's experts to MXFP4 operands before the matrix multiply.

\paragraph{Decode cost and speed.} Like EXL3 and QTIP~\citep{exllamav3,tseng2024qtip}, GRAFT decodes on the
general-purpose cores rather than on the tensor cores. For each weight, the decoder extracts the 16-bit state (a shift
and a mask), computes the hash (an integer multiply and a byte-sum dot product) and applies the value map (a fused
multiply-add and half an FP4 conversion). On Blackwell (B200) this takes about \final{K-B200-OFF} instructions in an
unflagged block and \final{K-B200-ON} with a flag, against \final{K-B200-OCT} for the per-matrix octave codebook.
The smaller representation reduces weight traffic, but whether this decode work is hidden and the saving turns into
throughput depends on overlapping the decode with memory access. Where decode itself sets the speed, as on HBM GPUs such as B200, a throughput model puts GRAFT's expert
products at an estimated 2--3$\times$ the time of native MXFP4, so on these GPUs its benefit is capacity. The decode
cost cannot be traded away cheaply. On \final{DEC-N} DeepSeek matrices (\final{DEC-HI-N} for the bit-sliced
decoders), every cheaper state-to-value map we screened (top-bit and register tables, two weights per state, linear and
bit-sliced decoders) raised weight error by \final{DEC-RANGE}, whereas a random permutation of GRAFT's own table raised
it by \final{DEC-PERM}. The byte-sum hash spreads the successors of a state better than random draws do. Hardware support
for trellis decoding in the tensor-core operand path, like today's support for 2:4 sparsity metadata and MX block
scales~\citep{rouhani2023mx}, would remove this cost for EXL3 and QTIP as well as for GRAFT.

\paragraph{Conclusion.} An MXFP4 checkpoint is a discrete source, and a quantizer that keeps the grid instead of
rotating it away can go well below the Gaussian-source curve for which rotation-based, Gaussian-matched coders are
designed. GRAFT does this inside EXL3's
trellis by changing only the value map. The new map applies an E2M1 conversion to EXL3's hash, with a one-bit MX
scale per block and a power-of-two tail at the hash centre, so the bulk decodes to MXFP4 blocks again. On
DeepSeek-V4.1-Flash, GRAFT has \final{X-MEAN2}$\times$ EXL3's held-out expert error at 3 bits, and its MXFP4 blocks
run on the model's own MXFP4 kernels with a per-column output scale and a sparse tail correction added. GRAFT
therefore stores the experts in fewer than 4 bits per weight while the bulk of the computation stays in the format the
model was trained in.
